\documentclass[12pt]{spieman}  % 12pt font required by SPIE;
%\documentclass[a4paper,12pt]{spieman}  % use this instead for A4 paper
\usepackage{amsmath,amsfonts,amssymb}
\usepackage{graphicx}
\usepackage{setspace}
\usepackage{tocloft}
\usepackage{lineno}
\linenumbers
\title{Wind-driven halo reconstruction and removal with wavefront sensor telemetry}

\author[a,*]{Author Name}
\affil[a]{Affiliation}

\renewcommand{\cftdotsep}{\cftnodots}
\cftpagenumbersoff{figure}
\cftpagenumbersoff{table}
\begin{document}
\maketitle

\begin{abstract}
The wind-driven halo (WDH) is a persistant noise artifact that arises due to the servo-lag inherent to all adaptive optics (AO) instruments. The standard procedure for dealing with this type of noise in science images is spatial
filtering. However, filtering out the WDH while simultaneously preserving signal from an extended object of interest is exceptionally challenging. Additionally, since the WDH changes in intensity and position angle through an observation, PCA-based algorithms (e.g., KLIP) that are commonly used to subtract the starlight need to be overly-aggressive to remove both the static and dynamic noise components. Since wavefront sensors (WFSs) continuously track the closed-loop residual wavefront error, WFS telemetry presents the ideal resource for combatting this type of noise artifact through postprocessing. Using archival WFS telemetry from MagAO-X, which is the extreme AO instrument for the 6.5-meter Magellan Clay telescope, we present a novel software tool for WDH reconstruction and removal. MagAO-X is equipped with a pyramid WFS capable of running high-order loop speeds of up to 3.7kHz and all of the WFS images are saved during data acquisiton. Given this, we detail how our new tool, windsoCC, cross-correlates the recorded closed-loop wavefront to measure the wind speed, direction, and relative strength of several individual layers of the atmosphere above Las Campanas Observatory. We then make use of the wind parameters learned through windsoCC to reconstruct and remove the WDH in individual MagAO-X science images by leveraging a parametric model. Notably, windsoCC facilitates a dramatic increase in the signal-to-noise at small inner-working angles in images of circumstellar disks and will benefit future efforts toward predictive control.
\end{abstract}

% Include a list of up to six keywords after the abstract
\keywords{adaptive optics, wavefront sensing, wind-driven halo, high-contrast imaging, post-processing, MagAO-X}

% Include email contact information for corresponding author
{\noindent \footnotesize\textbf{*}Author name, \linkable{email@example.edu} }

\begin{spacing}{2}   % use double spacing for rest of manuscript

\section{Introduction}
\label{sect:intro}

\begin{figure}
\begin{center}
\begin{tabular}{c}
\includegraphics[width=\linewidth]{camsci1_i_20230309_10_backend_file_mcmc_BestModel_Plot.jpg}
\end{tabular}
\end{center}
\caption
{ \label{fig:wdh-model}
Placeholder caption for the tentative WDH model figure. }
\end{figure}

\section{Observations and Data Reduction}
\label{sect:observations}

\section{WindsoCC pipeline description}
\label{sect:pipeline}

\subsection{Data partitioning}
\label{sect:data-partitioning}

\subsection{WFS data reduction}
\label{sect:wfs-data-reduction}

\subsection{Cross-correlating WFS data}
\label{sect:cross-correlating-wfs-data}

\subsection{Combining cross-correlation response maps}
\label{sect:combining-cross-correlation-response-maps}

\subsection{Tracking and measurement of wind}
\label{sect:tracking-and-measurement-of-wind}

\subsection{Quantifying wind attributes}
\label{sect:quantifying-wind-attributes}

\section{Forward modeling the WDH}
\label{sect:forward-modeling}

\subsection{Analytical model description}
\label{sect:analytical-model-description}

\subsection{Forward-modeling through KLIP}
\label{sect:forward-modeling-through-klip}

\subsection{Parameter estimation}
\label{sect:parameter-estimation}

\section{On-Sky demonstration of WDH correction}
\label{sect:on-sky-demonstration}

\subsection{Retrieved wind behavior}
\label{sect:retrieved-wind-behavior}

\subsection{WDH removal}
\label{sect:wdh-removal}

\section{Discussion}
\label{sect:discussion}

\section{Summary}
\label{sect:summary}

% \disclosures
\subsection*{Disclosures}


\subsection*{Code, Data, and Materials Availability}


\subsection*{Acknowledgments}


%%%%% References %%%%%

\bibliography{report}   % bibliography data in report.bib
\bibliographystyle{spiejour}   % makes bibtex use spiejour.bst

\listoffigures
\listoftables

\end{spacing}
\end{document}